%% kapa: exam=1 round=24 year=2013 period=2 form=A subject=appraisal_law subject_ko=부동산관계법규 seq=38 number=78
%% source: Q-Net 감정평가사 기출문제 (cst00309) artlSeq=5073802 file=extracted/1008/5073802/2013년 제24회 감정평가사 1차시험 기출문제/2교시(A형).hwp
%% conversion: kordoc 4.24.51 -> tools/gen1.py (auto); review=proofread-sonnet; flags=none
\begin{question}{78}
부동산등기법상 권리에 관한 등기 절차에 관한 설명으로 옳지 않은 것은?
\begin{choices}
  \item 등기원인에 권리의 소멸에 관한 약정이 있을 경우 신청인은 그 약정에 관한 등기를 신청할 수 있다.
  \item 등기명의인인 사람의 사망으로 권리가 소멸한다는 약정이 등기되어 있는 경우에 사람의 사망으로 그 권리가 소멸하였을 때에는, 등기권리자는 그 사실을 증명하여 단독으로 해당 등기의 말소를 신청할 수 있다.
  \item 등기권리자가 등기의무자의 소재불명으로 인하여 공동으로 등기의 말소를 신청할 수 없을 때에는 「민사소송법」에 따라 공시최고를 신청할 수 있다.
  \item 말소된 등기의 회복을 신청하는 경우에 등기상 이해관계 있는 제3자가 있을 때에는 그 제3자의 승낙이 있어야 한다.
  \item 대지권을 등기한 후에 한 건물의 권리에 관한 등기는, 그 등기에 건물만에 관한 것이라는 뜻이 부기되어 있을 때에도 대지권에 대하여 동일한 등기로서 효력이 있다.
\end{choices}
\end{question}
\answer{5}
